\title{跳表（Skip Lists）的应用与实现}
\author{叶家炜}
\date{Apr 19, 2026}
\maketitle
数据结构在高效查询中的重要性不言而喻，尤其是在处理大规模有序数据时。传统的有序链表如单链表，其查询时间复杂度为 O(n)，这在数据量增大时会成为明显的性能瓶颈。二叉搜索树虽然将平均查询复杂度优化到 O(log n)，但其平衡性依赖于复杂的旋转操作，且在并发场景下容易引入锁竞争。跳表作为一种随机化的平衡多层链表结构，应运而生，它通过巧妙的多层索引设计，在保持实现简单性的同时，实现了高效的查询性能。\par
跳表的显著优势在于其平均时间复杂度均为 O(log n)的查找、插入和删除操作。与红黑树或 AVL 树相比，跳表的实现更为简洁，空间开销也更小，通常只需额外约 33\%{} 的指针存储。更重要的是，跳表天生支持并发操作，其无锁设计使其在多线程环境中表现出色，避免了传统树结构常见的锁粒度问题。\par
本文将从跳表的基本原理入手，逐步深入其详细实现、应用场景、实际代码以及进阶主题。通过理论分析与伪代码解读，帮助读者全面理解这一高效数据结构，并提供实践指导。文章结构清晰，先原理后实现，再到应用与优化，最后探讨进阶扩展。\par
\chapter{2. 跳表的基本原理}
跳表的核心是一个多层索引结构，最底层是一个有序的单链表，而上层则是稀疏的“快递道”，即跳跃指针。这些指针允许我们在查找时快速跳过大量节点，从而加速搜索过程。每个节点不仅存储数据值，还包含一个前向指针数组，用于连接不同层级的下一个节点，以及当前节点的最大层级信息。层高的生成采用随机方式，通常基于几何分布，概率参数 p 取 0.5，这确保了结构的平衡性。\par
查找操作从跳表的顶层开始，从头节点出发，向右跳跃到最后一个小于目标值的节点，然后下降一层重复此过程，直至底层确认位置。这种逐层下降的策略，使得平均查找时间为 O(log n)。插入操作首先执行查找以定位插入点，同时记录每层的“前驱”节点，随后为新节点随机生成层高，并更新相应指针。删除则类似，找到前后节点后直接绕过目标节点更新指针。这些操作的平均时间复杂度均为 O(log n)，得益于随机层高的统计特性。\par
随机层高的数学基础源于几何分布。假设每个节点向上扩展到下一层的概率为 p=0.5，则节点 i 的层高服从几何分布，其期望值为\${}\${}\textbackslash{}frac\{{}1\}{}\{{}1-p\}{}=2\${}\${}。对于 n 个节点，整个跳表的期望最大层高为\${}\${}\textbackslash{}log\_{}\{{}1/p\}{} n\${}\${}，约为\${}\${}\textbackslash{}log\_{}2 n\${}\${}。这种随机化避免了结构退化成单链表的风险，因为高层的节点数量期望上呈指数衰减：第 k 层的节点数约为\${}\${}n \textbackslash{}cdot p\^{}k\${}\${}，从而保证了跳跃效率。\par
\chapter{3. 跳表的详细实现（伪代码 + 示例）}
跳表的核心数据结构可以用 C++ 风格伪代码定义如下。Node 结构体包含 value 成员存储数据，forward 数组存储最多 MAX\_{}LEVEL 层的前向指针，level 记录该节点的最大层级。SkipList 类维护头哨兵节点 head、最大层数 maxLevel 以及层高生成概率 probability，默认 0.5。这个设计简洁高效，空间开销主要来自指针数组，通常 MAX\_{}LEVEL 设为 16 或 32，足以应对亿级数据。\par
\begin{lstlisting}
struct Node {
    int value;
    Node* forward[MAX_LEVEL];  // 前向指针数组
    int level;                 // 当前节点最大层级
};
class SkipList {
    Node* head;                // 表头哨兵节点
    int maxLevel;              // 最大层数
    float probability;         // 层高生成概率 (默认 0.5)
};
\end{lstlisting}
随机生成层高的函数是跳表随机化的关键。这个函数从 level=1 开始，循环检查 rand() < MAX\_{}RAND * probability 的条件（MAX\_{}RAND 通常为 RAND\_{}MAX），若满足则 level 递增，直至达到 maxLevel 上限或随机失败。解读此代码：rand()产生 0 到 RAND\_{}MAX 的均匀随机数，乘以 probability 后与 rand()比较，等价于以概率 p 生成更高层。这种几何分布确保低层节点密集、高层稀疏，平均层高为 2，防止了最坏情况。\par
\begin{lstlisting}
int randomLevel() {
    int level = 1;
    while (rand() < MAX_RAND * probability && level < maxLevel)
        level++;
    return level;
}
\end{lstlisting}
查找操作 search 的关键在于高效定位。从当前层 level = maxLevel 开始，设置当前节点 cur = head。从头节点出发，在当前层向右跳跃：while 循环中，若 cur->forward[level]存在且其 value < target，则跳到该节点；否则 level--下降一层。到达底层后，若 cur->forward[0]正好等于 target 则找到，否则未命中。此过程时间复杂度 O(log n)，因为每层期望跳跃步数为 1/p=2。\par
插入操作 insert 首先调用 search 记录每层的“前驱”节点到 update 数组中，这些前驱是插入位置的左侧节点。随后创建新 Node，value 设为 key，随机调用 randomLevel()生成其 level，并初始化 forward 指针为 nullptr。然后，从新节点的 level 逐层更新：对于每一层 i，新节点的 forward[i]指向 update[i]->forward[i]，并将 update[i]->forward[i]指向新节点。此举确保指针横向和纵向的一致性，避免循环引用。\par
删除操作 erase 类似 search，先记录前后节点到 update 数组。对于目标节点的所有层级 i，从 update[i]->forward[i]直接跳到该节点的 forward[i]，绕过目标节点。最后释放内存。此操作不改变其他节点的层高，保持随机性不变。整体而言，这些操作的空间复杂度为 O(1)额外空间（update 数组大小为 maxLevel），时间为 O(log n)。\par
\chapter{4. 跳表的应用场景}
在实际生产环境中，跳表广泛应用于分布式数据库的索引层。例如 LevelDB 和 RocksDB 的 MemTable 使用跳表实现内存有序映射，支持高效的范围查询和并发读写。其无锁特性特别适合高吞吐场景，避免了树结构的重平衡锁。\par
Redis 的有序集合 ZSet 底层正是跳表实现，zskiplist.c 文件中可见其完整代码。这使得 ZSet 支持 O(log n)的范围查询如 ZRANGE，同时插入删除高效。日志系统如 Apache Kafka 的部分索引也借鉴跳表思想，实现追加写入加快速扫描。实时搜索引擎 Elasticsearch 在某些动态索引场景中采用类似结构，以应对频繁更新。\par
与其他数据结构对比，跳表的实现复杂度远低于红黑树或 AVL 树，后者需维护严格的平衡不变量。空间开销上，跳表每节点期望 1/(1-p)=2 个指针，总空间约 1.33n，非常经济。并发支持是其杀手锏，无需锁即可实现线性化语义，而树结构往往需读写锁。范围查询时，跳表从定位点顺序遍历底层链表，效率与 B+ 树相当，但后者空间更高且实现复杂。\par
跳表并非完美，其随机化导致理论最坏 O(n)，虽概率指数级小。指针跳转也对 CPU 缓存不友好。优化策略包括固定层高以提升预测性，或预热缓存以减少 miss。\par
\chapter{5. 实际代码实现与测试}
完整 C++ 实现需定义 MAX\_{}LEVEL=16，头节点初始化所有 forward 为 nullptr，maxLevel=1。insert 函数中处理重复 key：若 search 找到则更新 value 或忽略。search 返回 pair<Node*, int>，Node 为最近前驱，int 为层级。delete 需两次 search 确认前后，并处理空表边界。Python 实现类似，用列表模拟 forward 数组，random.random() < p 生成层高；Java 用 AtomicReferenceArray 支持 CAS 无锁。\par
性能测试中，与 std::set 对比插入 100 万随机 int，跳表查询延迟通常快 20\%{}-50\%{}，因无旋转开销。基准代码可用 Google Benchmark 框架，测量 QPS 和 P99 延迟。常见问题如指针循环可用 DFS 检测，内存泄漏用 valgrind 追踪。多线程安全版加读写锁，或用 CAS 实现无锁。\par
\chapter{6. 进阶主题}
无锁并发跳表使用 CAS 操作指针更新。插入时，先用乐观定位记录 update，再 CAS 尝试链接新节点，若失败重试。此设计如 Harris 链表融合，确保 ABA 问题通过标记位解决。\par
持久化跳表针对 NVM 优化，节点用原子写持久化指针，结合 CLWB 指令保证顺序。分区跳表引入 sharding，将键空间分片到多跳表，提升并行度，如 TiDB 的索引层应用。\par
\chapter{7. 结论}
跳表的核心价值在于简单高效的概率平衡，用随机化换取确定性性能。学习它深化了对随机数据结构的理解。从 Redis 源码入手，实现基准测试，是最佳实践路径。\par
\chapter{8. 参考资料与扩展阅读}
原论文 William Pugh 的《Skip Lists: A Probabilistic Alternative to Balanced Trees》（1990）奠定基础。Redis zskiplist.c 和 LevelDB 源码提供实战参考。《Redis 设计与实现》和《数据库系统概念》有深入讨论。工具如 perf 和 valgrind 助性能调优。\par
